%% toc: P15 159 159.7,159.10
\chapter{A Probability Space Without a Regular Conditional Probability}
\label{sec:p15-w02c-p15-008-no-rcp}

\Cref{def:p15-regular-conditional-law} says that a regular conditional
probability on the original outcome space, the case \(S=\Omega\),
\(\mathcal S=\mathcal F_0\) with \(X\) the identity, needs an existence
hypothesis on that space. This chapter proves the claim by constructing a
probability space and a sub-sigma-algebra for which no regular conditional
probability exists, in either the sense of
\cref{def:p15-regular-conditional-law} or Durrett's weaker sense of
\cref{def:p15-w02b-p15-008-rcd}. Durrett states that regular conditional
distributions do not always exist and credits the first example to
Dieudonn\'e (1948), referring to Doob (1953), p.~624, and Faden (1985)
for later work, without giving an example
\cite[p.~214]{DurrettPTEAuthor2019}. None of those works is held for
this handbook. The model below is constructed and proved here from the
measure theory of Part~P15 and one use of the axiom of choice; it is not
offered as a reconstruction of any of those examples.

\emph{Inputs.} Lebesgue measure on the Borel sets of \(\mathbb R\)
(\cref{thm:p15-lebesgue-unit-interval}) with its construction from
interval covers (\cref{thm:p15-premeasure-extension,thm:p15-stieltjes-construction}),
the uniform probability space \((\Omega_U,\mathcal F_U,\mathbb U)\), and
the existence and almost-everywhere uniqueness of conditional expectations
(\cref{thm:p15-conditional-existence}).
\emph{Conclusions.} A probability \(P_H\) on a sigma-algebra
\(\mathcal F_H\supseteq\mathcal F_U\) of subsets of \((0,1)\) such that,
with \(\mathcal G=\mathcal F_U\), no function
\(\mu:(0,1)\times\mathcal F_H\to[0,1]\) is a regular conditional
probability. The obstruction lies in the target sigma-algebra: real random
variables on the same space, and the identity map into
\(((0,1),\mathcal F_U)\), all have regular conditional laws given
\(\mathcal F_U\).

\section{A set without Borel subsets of positive length}
\label{sec:p15-w02c-p15-008-thick-set}

\begin{hbconvention}{def:p15-w02c-p15-008-groups}{The groups of the construction and the use of choice}
Put
\[
 \begin{gathered}
 L\defeq\{m+n\sqrt2:m,n\in\mathbb Z\},\\
 L_0\defeq\{2m+n\sqrt2:m,n\in\mathbb Z\},\\
 L_1\defeq\{2m+1+n\sqrt2:m,n\in\mathbb Z\}.
 \end{gathered}
\]
For \(A,B\subseteq\mathbb R\) and real \(t\) write
\(A+B=\{a+b:a\in A,\ b\in B\}\), \(A-B=\{a-b:a\in A,\ b\in B\}\) and
\(A+t=A+\{t\}\). The construction of \cref{thm:p15-w02c-p15-008-thick-set}
uses the axiom of choice once, to select one point from each member of
a family of disjoint nonempty sets. That is the only use of choice in
this chapter; every later result depends on it through the set \(H\).
\end{hbconvention}

\begin{hblemma}{thm:p15-w02c-p15-008-groups}{proved=annex:p15-w02c-p15-008-groups}{The even and odd cosets}
\begin{enumerate}
\item \(L\) and \(L_0\) are additive subgroups of \(\mathbb R\),
      \(L_1=L_0+1\), and \(L\) is the disjoint union of \(L_0\) and
      \(L_1\). Moreover \(L_1+L_0=L_1\).
\item For every \(\delta>0\) the set \(L_1\cap(-\delta,\delta)\) is
      nonempty.
\end{enumerate}
\end{hblemma}

\begin{hblemma}{thm:p15-w02c-p15-008-steinhaus}{proved=annex:p15-w02c-p15-008-steinhaus}{Differences of a set of positive length}
Let \(F\subseteq\mathbb R\) be Borel with \(0<\lambda(F)<\infty\).
Then there is \(\delta>0\) such that \((-\delta,\delta)\subseteq F-F\).
\end{hblemma}

\begin{hbtheorem}{thm:p15-w02c-p15-008-thick-set}{proved=annex:p15-w02c-p15-008-thick-set}{A set splitting Borel sets of positive length}
Choose \(V\subseteq\mathbb R\) containing exactly one point of each coset
\(x+L\), \(x\in\mathbb R\), and put
\[
 H\defeq V+L_0 .
\]
Then \(\mathbb R\setminus H=V+L_1=H+1\), and every Borel set
\(F\subseteq H\) and every Borel set \(F\subseteq\mathbb R\setminus H\)
has \(\lambda(F)=0\). Consequently, with
\(H_U\defeq H\cap(0,1)\) and \(H_U^c\defeq(0,1)\setminus H_U\), every
member of \(\mathcal F_U\) contained in \(H_U\) or in \(H_U^c\) has
\(\mathbb U\)-measure zero, and \(H_U\notin\mathcal F_U\).
\end{hbtheorem}

Put differently, every Borel set \(B\subseteq(0,1)\) of positive length
meets both \(H_U\) and \(H_U^c\): otherwise \(B\) would lie inside one of
them.

\section{The model}
\label{sec:p15-w02c-p15-008-model}

\begin{hbdefinition}{def:p15-w02c-p15-008-model}{Adjoining the set to the Borel sets}
With \(H_U\) and \(H_U^c\) as in \cref{thm:p15-w02c-p15-008-thick-set},
let
\[
 \mathcal F_H\defeq\bigl\{(B_1\cap H_U)\cup(B_2\cap H_U^c):
                    B_1,B_2\in\mathcal F_U\bigr\}
\]
and, for \(D=(B_1\cap H_U)\cup(B_2\cap H_U^c)\in\mathcal F_H\),
\[
 P_H(D)\defeq\tfrac12\bigl(\mathbb U(B_1)+\mathbb U(B_2)\bigr).
\]
The outcome space is \(\Omega_U=(0,1)\), and the conditioning
sigma-algebra is \(\mathcal G\defeq\mathcal F_U\).
\end{hbdefinition}

\begin{hbtheorem}{thm:p15-w02c-p15-008-model}{proved=annex:p15-w02c-p15-008-model}{The enlarged probability space}
\begin{enumerate}
\item \(\mathcal F_H\) is a sigma-algebra on \((0,1)\) that contains
      \(\mathcal F_U\) and \(H_U\).
\item \(P_H\) is well defined: its value does not depend on the
      representation of \(D\). It is a probability measure on
      \(\mathcal F_H\), and \(P_H(B)=\mathbb U(B)\) and
      \(P_H(B\cap H_U)=P_H(B\cap H_U^c)=\tfrac12\mathbb U(B)\) for every
      \(B\in\mathcal F_U\).
\item In \((\Omega_U,\mathcal F_H,P_H)\), for every \(B\in\mathcal F_U\)
      the indicator \(\mathbf1_B\) is a version of
      \(P_H(B\mid\mathcal F_U)\), and the constant \(\tfrac12\) is a
      version of \(P_H(H_U\mid\mathcal F_U)\). Every version of
      \(P_H(H_U\mid\mathcal F_U)\) equals \(\tfrac12\) outside an
      \(\mathcal F_U\)-measurable \(P_H\)-null set.
\end{enumerate}
\end{hbtheorem}

\section{Nonexistence of a regular conditional probability}
\label{sec:p15-w02c-p15-008-nonexistence}

\begin{hbtheorem}{thm:p15-w02c-p15-008-no-rcp}{proved=annex:p15-w02c-p15-008-no-rcp}{No regular conditional probability given the Borel sets}
In the probability space \((\Omega_U,\mathcal F_H,P_H)\) of
\cref{def:p15-w02c-p15-008-model}, let \(X\) be the identity map of
\((\Omega_U,\mathcal F_H)\) onto itself and \(\mathcal G=\mathcal F_U\).
\begin{enumerate}
\item No function \(\mu:\Omega_U\times\mathcal F_H\to[0,1]\) is a regular
      conditional distribution for \(X\) given \(\mathcal G\) in the sense
      of \cref{def:p15-w02b-p15-008-rcd}. That is, the space has no regular
      conditional probability given \(\mathcal F_U\) in Durrett's sense.
\item No probability kernel from \((\Omega_U,\mathcal F_U)\) to
      \((\Omega_U,\mathcal F_H)\) is a regular conditional law of \(X\)
      given \(\mathcal G\) in the sense of
      \cref{def:p15-regular-conditional-law}.
\item The measurable space \((\Omega_U,\mathcal F_H)\) admits no
      measurable isomorphism onto a Borel subset of \(\mathbb R\) of the
      kind assumed in \cref{thm:p15-standard-conditional-law}.
\item The point-mass kernel \(\delta(\omega,B)=\mathbf1_B(\omega)\),
      \(\omega\in\Omega_U\), \(B\in\mathcal F_U\), is a regular
      conditional law, given \(\mathcal G\), of the identity map
      \((\Omega_U,\mathcal F_H)\to(\Omega_U,\mathcal F_U)\). Every real
      random variable on \((\Omega_U,\mathcal F_H,P_H)\) has a regular
      conditional law given \(\mathcal G\).
\end{enumerate}
\end{hbtheorem}

The proof of part~(i) locates the obstruction. Countably many of the
required identities, one for each interval \((0,q]\) with rational
\(q\), force \(\mu(\omega,\cdot)\) to give mass one to the singleton
\(\{\omega\}\) for almost every \(\omega\). A measure with that property
assigns \(H_U\) the value \(\mathbf1_{H_U}(\omega)\), which is \(0\) or
\(1\). The identity for the single further event \(H_U\) requires the
value \(\tfrac12\) almost everywhere. Since the exceptional sets involved
are countably many null sets, some outcome lies outside all of them, and
the two requirements contradict each other there.

\begin{hbexample}{def:p15-w02c-p15-008-worked}{Two natural candidates and where each fails}
Take \(B=(0,\tfrac12]\in\mathcal F_U\). By part~(ii) of
\cref{thm:p15-w02c-p15-008-model},
\[
 P_H(B)=\tfrac12,\qquad P_H(B\cap H_U)=\tfrac14,\qquad
 P_H(H_U)=\tfrac12 .
\]
\begin{enumerate}
\item \emph{Point masses.} Put \(\mu_1(\omega,A)=\mathbf1_A(\omega)\) for
      \(A\in\mathcal F_H\). Each \(\mu_1(\omega,\cdot)\) is a probability
      measure, and for \(A\in\mathcal F_U\) the function
      \(\mu_1(\cdot,A)=\mathbf1_A\) is a version of
      \(P_H(A\mid\mathcal F_U)\); for instance \(\mu_1(\cdot,B)\) is
      \(1\) on \((0,\tfrac12]\) and \(0\) elsewhere. For \(A=H_U\) the
      function \(\mathbf1_{H_U}\) is not \(\mathcal F_U\)-measurable,
      because \(H_U\notin\mathcal F_U\) by
      \cref{thm:p15-w02c-p15-008-thick-set}, so it is not a version of
      \(P_H(H_U\mid\mathcal F_U)\). Condition~(i) of
      \cref{def:p15-w02b-p15-008-rcd} fails at the event \(H_U\).
\item \emph{The unconditional law.} Put \(\mu_2(\omega,A)=P_H(A)\). This is
      a probability kernel, and \(\mu_2(\cdot,H_U)=\tfrac12\) is a version
      of \(P_H(H_U\mid\mathcal F_U)\). But \(\mu_2(\cdot,B)=\tfrac12\),
      while every version of \(P_H(B\mid\mathcal F_U)\) equals
      \(\mathbf1_B\) almost everywhere, and \(\tfrac12\neq\mathbf1_B(\omega)\)
      for every \(\omega\). Condition~(i) fails at the event \(B\), on a set
      of probability one.
\end{enumerate}
The first candidate is right on \(\mathcal F_U\) and wrong at \(H_U\); the
second is right at \(H_U\) and wrong on \(\mathcal F_U\).
\cref{thm:p15-w02c-p15-008-no-rcp} shows that no function is right on
both, even outside a null set of outcomes.
\end{hbexample}

\begin{hbwarning}{def:p15-w02c-p15-008-limits}{What the model shows and what it does not}
\begin{enumerate}
\item The model shows that existence of a regular conditional probability
      on \((\Omega,\mathcal F_0)\) is not automatic. It does not identify a
      necessary condition: \cref{thm:p15-standard-conditional-law} gives a
      sufficient one, and part~(iii) of
      \cref{thm:p15-w02c-p15-008-no-rcp} only shows that this model
      violates it.
\item The construction of \(H\) uses the axiom of choice. Nothing here
      decides whether a model of this kind can be given without it.
\item The probability space itself is not the source of the failure.
      Part~(iv) of \cref{thm:p15-w02c-p15-008-no-rcp} gives regular laws
      for every real random variable and for the identity into
      \((\Omega_U,\mathcal F_U)\); the failure appears only when the target
      sigma-algebra is enlarged by the one event \(H_U\).
\item The model is not a sequence of exchangeable coordinates. It does not
      exhibit an exchangeable law that is not presentable in the sense of
      \cref{def:p15-w02b-p15-022-presentable}; part~(iv) of
      \cref{thm:p15-w02b-p15-022-presentable} concerns such laws.
\end{enumerate}
\end{hbwarning}

\section{Relationships}
\label{sec:p15-w02c-p15-008-relationships}

\begin{itemize}
\item \Cref{thm:p15-w02c-p15-008-no-rcp} establishes the sentence in
      \cref{def:p15-regular-conditional-law} that the special case
      \(S=\Omega\), \(\mathcal S=\mathcal F_0\) needs an existence
      hypothesis, and it supplies the demonstration behind item~(iii) of
      \cref{def:p15-w02b-p15-008-limits}.
\item The countable-cocountable model of
      \cref{thm:p15-w02b-p15-008-ae-counterexample} has a regular
      conditional law; its failure concerns the repair of one given
      source-sense law. In the present model no law of either kind exists.
\item Durrett's Theorem~4.1.17 states that regular conditional
      distributions exist when the target space is nice
      \cite[p.~215]{DurrettPTEAuthor2019}; the handbook form is
      \cref{thm:p15-standard-conditional-law}, which part~(iii) of
      \cref{thm:p15-w02c-p15-008-no-rcp} shows cannot apply to
      \((\Omega_U,\mathcal F_H)\).
\item Durrett explains the difficulty by exceptional null sets that
      accumulate over many countable disjoint families
      \cite[p.~214]{DurrettPTEAuthor2019}; that mechanism is set out in
      \cref{def:p15-w02c-p15-002-rcd-nonexistence} with
      \cref{thm:p15-w02c-p15-002-ae-additivity}. In the present model only
      countably many exceptional sets are used. The contradiction comes from
      the incompatibility, at a single outcome, of the values forced on
      \(\mathcal F_U\) with the value required at \(H_U\).
\end{itemize}

\emph{Sources.} Durrett, \emph{Probability: Theory and Examples},
version~5 dated 11 January 2019, p.~214 for the definition of a
regular conditional distribution and probability, the statement that such
objects do not always exist with its attribution to Dieudonn\'e (1948),
Doob (1953) and Faden (1985), and the accumulation explanation; p.~215 for
Theorem~4.1.17 \cite[pp.~214--215]{DurrettPTEAuthor2019}. The groups
\(L_0,L_1\), the set \(H\), the difference lemma and the model
\((\Omega_U,\mathcal F_H,P_H)\) are written here without a source
attribution.
